% META: {"id": "TSPE-GEOMETRIE-ESPACE-RE-C14", "chapitre": "TSPE-GEOMETRIE-ESPACE", "type_objet": "remediation", "capacites_codes": ["C14"], "status": "generated"}

%---------------------------------------------------------------
% FICHE DE REMEDIATION — C14 : coordonnees du projete orthogonal
% Obstacle : annuler une coordonnee quel que soit le plan
%---------------------------------------------------------------

\subsection*{Erreur fréquente}

\erreurFrequente{
Projeter un point sur un plan en annulant sa troisième coordonnée.
}

\textbf{Pourquoi c'est faux.} Annuler la cote ne projette que sur le plan
$z=0$. Pour un plan quelconque, la direction de projection est celle du
vecteur normal, qui n'a aucune raison d'être verticale.

\contreexemple{
Sur le plan $x+y+z-3=0$, le point $M(3;3;3)$ ne se projette pas en
$(3;3;0)$~: ce point ne vérifie pas $3+3+0-3=0$. Il n'est donc même pas dans
le plan.
}

\subsection*{À retenir}

Le projeté orthogonal $H$ de $M$ sur $\mathcal{P}$ s'obtient en suivant la
normale~: on paramètre la droite passant par $M$ et dirigée par $\vec{n}$,
puis on cherche son point d'intersection avec le plan.

\subsection*{Méthode corrective}

\begin{enumerate}
  \item Écrire la droite $\Delta$ passant par $M$, de vecteur directeur
        $\vec{n}$~: $X=M+t\,\vec{n}$.
  \item Reporter ces coordonnées dans l'équation du plan.
  \item Résoudre en $t$~: la solution est unique.
  \item Reporter ce $t$ dans $X$ pour obtenir $H$, puis vérifier que $H$
        appartient bien au plan.
\end{enumerate}

%---------------------------------------------------------------

\begin{exercice}{TSPE-GEOESPACE-RE-C14-EX1}{2}{14}
Repère orthonormé. Soit $\mathcal{P}:\ x+y+z-3=0$ et $M(3;3;3)$.
\begin{enumerate}
  \item Vérifier que $(3;3;0)$ n'appartient pas à $\mathcal{P}$.
  \item Écrire une représentation paramétrique de la droite passant par $M$
        et orthogonale à $\mathcal{P}$.
  \item Déterminer les coordonnées du projeté orthogonal $H$ de $M$.
  \item Calculer $MH$ et comparer à $d(M,\mathcal{P})$.
\end{enumerate}
\end{exercice}

% BEGIN-VERIFY
% from sympy import Matrix, symbols, solve, Eq, sqrt, simplify
% t = symbols('t')
% n = Matrix([1,1,1]); d = -3
% M = Matrix([3,3,3])
% # 1. (3,3,0) n'est pas dans le plan
% assert n.dot(Matrix([3,3,0])) + d != 0
% # 3. intersection de M + t n avec le plan
% sol = solve(Eq(n.dot(M + t*n) + d, 0), t)
% assert len(sol) == 1
% H = M + sol[0]*n
% assert list(H) == [1, 1, 1]
% assert n.dot(H) + d == 0
% # 4. MH egale la distance point-plan
% assert simplify((H - M).norm() - sqrt(12)) == 0
% assert simplify((H - M).norm() - abs(n.dot(M) + d)/n.norm()) == 0
% END-VERIFY

\begin{corrige}{TSPE-GEOESPACE-RE-C14-EX1}
\begin{enumerate}
  \item $3+3+0-3=3\neq0$~: le point $(3;3;0)$ n'est pas dans $\mathcal{P}$.
        Annuler la cote ne projette donc pas sur ce plan.
  \item $\vec{n}(1;1;1)$ est normal à $\mathcal{P}$, d'où
        $x=3+t$, $y=3+t$, $z=3+t$.
  \item En reportant~: $(3+t)+(3+t)+(3+t)-3=0$, soit $3t+6=0$ et $t=-2$.
        Donc $H(1;1;1)$, et l'on vérifie $1+1+1-3=0$.
  \item $MH=\sqrt{(-2)^2+(-2)^2+(-2)^2}=\sqrt{12}=2\sqrt{3}$. Par la formule,
        $d(M,\mathcal{P})=\dfrac{\lvert3+3+3-3\rvert}{\sqrt{3}}
        =\dfrac{6}{\sqrt{3}}=2\sqrt{3}$~: les deux coïncident, comme attendu.
\end{enumerate}
\end{corrige}
